% §8 Conclusion
Real-time responsiveness in cascade voice agents is two-sided, and this paper supplies the side
that has neither a standard definition nor a bound: the cessation floor.
We close the gap with a two-floor model and a phase-conditioned result: cessation grows with
in-flight audio up to the committed buffer depth, generation cancellation is prompt and
independent of the delivery tail, and resumption after interruption is dominated by tool
round-trip cost unless tools run on a background clock.

The enabling mechanism is orchestration, not any single component.
Decoupling generation, delivery, and tool clocks makes an agent generation-interruptible and
execution-persistent at once, and an API-serializability lemma guarantees that cancelling
generation never strands a pending tool batch.
A small, tested orchestrator verifies the semantics, and a monotonic-clock smoke anchor
instantiates the quantitative shape: cessation linear in in-flight audio, prompt
depth-independent cancellation, and tool-dominated resumption.

The practical conclusion answers the debate that motivated this work.
For an assistant whose turns are short and memory-bound, orchestration choice is noise under a
fixed in-flight policy.
For an assistant whose turns routinely call tools -- a bank, a card service, a multi-step
skill -- an idiomatically decoupled runtime avoids paying each round trip inside the
user's perceived silence and manages the interruption tail.
Framework arguments settle no data; the two floors settle the question: onset and cessation
together describe perceived responsiveness, and the onset guarantee now extends to the
orchestration that schedules reasoning, tools, and interruptions.
